\section{Exact height and displacement atoms}\label{c:atoms-section}

We now turn strip regularity into estimates for a bridge with a prescribed
terminal port.  Throughout this section, $p$ is the probability law of one
irreducible bridge from a fixed port.  Write $H$, $Z$, $L$, and $D$
for its positive integral height, transverse displacement, number of visited
triangle centers, and diameter.  The vector $X=(H,Z)$ takes values on the
actual row-staggered lattice; an atom outside this lattice is zero.  For
independent copies, put
\[
 S_k=X_1+\cdots+X_k,\qquad
 T_k=L_1+\cdots+L_k.
\]
Independence is used between whole pieces, never between coordinates of
one piece.

We use the renewal identity and one-piece estimates stated earlier:
\begin{equation}\label{c:renewal-atom-inputs}
 \sum_{h\ge0}B_hz^h=\frac1{1-\mathbb E_p z^H},\qquad
 B_0=1,\quad B_h\asymp(1+h)^{-1/4},
\end{equation}
and, for $x,t\ge1$ and $u\downarrow0$,
\begin{equation}\label{c:tail-atom-inputs}
 \begin{split}
 p(H>x)+p(D>x)&\le Cx^{-3/4},\\
 p(L>t)&\le t^{-9/16+o(1)},\qquad
 1-\mathbb E_p e^{-uL}=u^{9/16+o(1)}.
 \end{split}
\end{equation}
The first identity implies
$1-\mathbb E_p e^{-sH}\asymp s^{3/4}$.
These are statements under $p$, before conditioning on an endpoint.

\subsection{An atom of the one-piece height}

The extra information supplied by strip regularity is one full power of
height decay beyond the tail bound.  This is the estimate that rules out
bridges made from too few irreducibles at one prescribed endpoint.

\begin{proposition}\label{c:height-atom}
For every $\eta>0$ there is $C_\eta<\infty$ such that
\[
 p(H=n)\le C_\eta n^{-7/4+\eta}\qquad(n\ge1).
\]
\end{proposition}

\begin{proof}
Let $b(z)=\sum_{m\ge0}B_mz^m$.  Fix $n\ge2$, set
$r=e^{-1/n}$, and write $s=\min\{1,n^{-1}+|\theta|\}$ for
$-\pi\le\theta\le\pi$.  The height tail gives
\[
 |1-\mathbb E_p(re^{i\theta})^H|
 \le C\mathbb E_p\min\{1,sH\}\le Cs^{3/4}.
\]
Thus $|b(re^{i\theta})|\ge cs^{-3/4}$.  In particular $b$ has no
zero on this circle, as also follows from its renewal representation.

For every fixed integer $k\ge0$ and every $\delta>0$, the strip
differences of Proposition~\ref{c:strip-smoothness} imply
\begin{equation}\label{c:b-angular-derivatives}
 \left|\partial_\theta^k b(re^{i\theta})\right|
 \le C_{k,\delta}s^{-3/4-k-\delta}.
\end{equation}
Here are the details of the summation estimate.  Choose a smooth dyadic
partition of unity on the positive integers.  On $m\asymp M$, discrete
product rules and the strip differences bound every $j$th difference of
$m^kB_mr^m$ times the cutoff by
$C_{j,k,\delta}M^{k-1/4-j+\delta}e^{-M/(Cn)}$.
Summation by parts $j$ times, together with the trivial sum bound, gives
\[
 C_{j,k,\delta}M^{3/4+k+\delta}
 e^{-M/(Cn)}(1+M|\theta|)^{-j}.
\]
Choose $j>3/4+k+\delta$ and sum over dyadic $M$.  The terms with
$M\le s^{-1}$ form a geometric sum; the remaining terms are bounded by
the oscillatory factor or the exponential cutoff.  This proves
\eqref{c:b-angular-derivatives}.

Differentiating the reciprocal twice now gives
\[
 \left|\partial_\theta^2 b(re^{i\theta})^{-1}\right|
 \le C_\delta s^{-5/4-2\delta}.
\]
Consequently its integral over $[-\pi,\pi]$ is at most
$C_\delta n^{1/4+2\delta}$.  Since
$\mathbb E_pz^H=1-b(z)^{-1}$, the $n$th Fourier coefficient on this
circle, integrated by parts twice, satisfies
\[
 r^n p(H=n)
 \le\frac1{2\pi n^2}
 \int_{-\pi}^{\pi}\left|
 \partial_\theta^2 b(re^{i\theta})^{-1}\right|\,d\theta
 \le C_\delta n^{-7/4+2\delta}.
\]
The boundary terms vanish by periodicity, $r^n=e^{-1}$, and $\delta$
is arbitrary.
\end{proof}

\subsection{Two forms of displacement smoothing}

We need both a planar atom bound and a bound summed over possible heights.
The latter is stronger than merely summing the planar bound over all
heights, since the height distribution has an infinite tail.

\begin{proposition}\label{c:planar-atoms}
There is $C<\infty$ such that, for every integer $k\ge1$,
\begin{equation}\label{c:two-atom-bounds}
 \sup_w\Pr(S_k=w)\le Ck^{-8/3},\qquad
 \sum_{m\in\mathbb Z}\sup_{H(w)=m}\Pr(S_k=w)
 \le Ck^{-4/3}.
\end{equation}
\end{proposition}

\begin{proof}
We first establish a directional window for one piece.  There are constants
$0<c_0<C_0<\infty$ and $c_1>0$ such that, uniformly over unit vectors $e$
and all sufficiently large $x$,
\begin{equation}\label{c:directional-window}
 p\{c_0x\le |e\cdot X|\le C_0x\}\ge c_1x^{-3/4}.
\end{equation}
Use bridges of heights in $[x,2x]$ and terminal ports in a transverse
interval of length comparable to $x$.  The point-to-point connector lower
bound is $cx^{-5/4}$ at each such port.  The two intervals of bounded
positive and negative slope may be chosen so that their displacements
have projection at least a fixed multiple of $x$ in any prescribed
direction.  Compactness of the unit circle makes the constants uniform.
Summing over heights and ports gives critical mass at least $cx^{3/4}$.

By the renewal factorization, this is a sum of probabilities for $S_k$.
The terms $k>Ax^{3/4}$ have total mass at most
\[
 \sum_{k>Ax^{3/4}}e^2
   (\mathbb E_p e^{-H/x})^k
 \le Cx^{3/4}e^{-cA}.
\]
Fix $A$ large.  Some $k\le Ax^{3/4}$ therefore has probability bounded
below of producing an absolute projection comparable to $x$.
The probability of an increment with $D>C_0x$ is at most
$CA C_0^{-3/4}$, while
\[
 \mathbb E\sum_{i\le k}|e\cdot X_i|
       \mathbf1_{\{|e\cdot X_i|<c_0x\}}
 \le CA c_0^{1/4}x.
\]
The second estimate follows by integrating the diameter tail; fixed
lattice conversion constants are absorbed in $C$.  Choose $C_0$ large
and $c_0$ small.  With a fixed positive probability an increment then
lies in the window in \eqref{c:directional-window}.  A union bound and
$k\le Ax^{3/4}$ prove that display.

Let $\varphi(\theta)=\mathbb E e^{i\theta\cdot X}$, using any fixed
integer coordinates for the displacement lattice.  An independent copy
$X'$ lies in a fixed bounded set with probability bounded below.  Apply
\eqref{c:directional-window} in direction $\theta/|\theta|$, with
$x$ a sufficiently small fixed multiple of $|\theta|^{-1}$.  On the
resulting event, $|\theta\cdot(X-X')|$ stays between two positive
constants less than $\pi$.  Hence
\[
 1-|\varphi(\theta)|^2
 =\mathbb E[1-\cos(\theta\cdot(X-X'))]
 \ge c|\theta|^{3/4}
\]
near zero.  The directional windows also show that the affine span of the
support is two-dimensional.  Its difference group is therefore a
full-rank lattice subgroup.  There are only finitely many points of the
lattice torus where $|\varphi|=1$; near each, its modulus is a translate
of the modulus near zero.  On the remaining compact set it is strictly
less than one.  Fourier inversion gives
\[
 \sup_w\Pr(S_k=w)
 \le C\int_0^1 e^{-ckr^{3/4}}r\,dr+Ce^{-ck}
 \le Ck^{-8/3}.
\]
This also explains why no aperiodicity assumption beyond the actual
support lattice is needed.

It remains to sum the heightwise supremums.  For a nonnegative measure
$q$ on the displacement lattice, set
\[
 \|q\|_{1,\infty}
 =\sum_m\sup_{H(w)=m}q(w).
\]
Convolution satisfies
$\|q*r\|_{1,\infty}\le\|q\|_{1,\infty}\|r\|_1$.
The first bound of \eqref{c:two-atom-bounds} immediately bounds the sum
over $m\le2k^{4/3}$ by $Ck^{-4/3}$.

For a dyadic shell $m\asymp\Lambda k^{4/3}$, $\Lambda\ge2$, split the $k$ pieces
into sixteen independent blocks with sizes comparable to $k$; bounded
$k$ can be absorbed into the constant.  Each block has planar atoms at
most $Ck^{-8/3}$.  Also a block height exceeds $t\ge k^{4/3}$ with
probability at most $Ckt^{-3/4}$: split at increments larger than $t$
and use the truncated first moment for the rest.

Suppose first that one block has height at most $k^{4/3}\sqrt\Lambda$.
For large $\Lambda$, another has height at least $c\Lambda k^{4/3}$.  The first
restricted block has $\|\cdot\|_{1,\infty}$ at most
$C\sqrt\Lambda\,k^{-4/3}$ by its planar atom bound, while the latter event
has probability at most $C\Lambda^{-3/4}$.  Convolve with the remaining blocks
and sum over the finitely many choices.  This gives
$C\Lambda^{-1/4}k^{-4/3}$ for the shell.  If all sixteen block heights exceed
$k^{4/3}\sqrt\Lambda$, impose the planar atom on one block and the probability
bound $C\Lambda^{-3/8}$ on each of the other fifteen.  There are
$O(\Lambda k^{4/3})$ heights in the shell, so this case costs at most
$C\Lambda^{-37/8}k^{-4/3}$.  Both bounds are summable over dyadic $\Lambda$.
Increasing the constant handles the finitely many small $\Lambda$ and bounded
$k$, completing the proof.
\end{proof}

\section{Prescribed ports on parallel sides}\label{c:parallel-section}

The strip estimates above concern a fixed initial port and all terminal
ports.  We now prescribe the second port as well.  The connector
construction gives the required normalization: when the height is $h$
and the transverse displacement is $O(h)$, the critical mass between the
two ports is at least $ch^{-5/4}$.  The constant is uniform over fixed
compact nondegenerate scaled geometries.  The same construction lies in
a convex pure-sided domain if both ports have clearance comparable to
$h$ from every side other than the opposite parallel sides containing
them.  We will show that the exceptional mass in the entire strip is
$o(h^{-5/4})$.

\begin{theorem}\label{c:parallel}
Fix a bound on the ratio of transverse displacement to height.  For
compatible prescribed ports on two parallel pure cuts at integral height
$h\to\infty$, normalize the critical bridge weights by their total
mass.  For every $\epsilon>0$,
\[
 \Pr\{h^{4/3-\epsilon}\le L
                       \le h^{4/3+\epsilon}\}\longrightarrow1.
\]
The convergence is uniform under the fixed displacement bound.  It also
holds in convex pure-sided domains in a fixed compact family of
nondegenerate scaled geometries, with the ports on opposite parallel
sides and with order-$h$ clearance from every other side.
\end{theorem}

\begin{proof}
It is enough to treat $0<\epsilon<1$.  Let $w$ be the prescribed
displacement, with $H(w)=h$.  Unique
factorization writes its bridge mass as
$\sum_{k\ge1}\Pr(S_k=w)$.  We first discard very small and very large
numbers of pieces without imposing a length restriction.

Fix a small $\xi>0$.  Put $K=h^{3/4-\xi}$ and
$a=h^{1-\tau}$, where $0<\tau<8\xi/21$.  For $k\le K$, the event
$\sum_{i\le k}H_i=h$ with every $H_i<a$ has probability at most
\[
 e^{-h/a}\bigl(\mathbb E[e^{H/a};H<a]\bigr)^k
 \le\exp\{-h^\tau+Cka^{-3/4}\}.
\]
We used $e^t-1\le et$ for $0\le t\le1$ and
$\mathbb E[H;H<a]\le Ca^{1/4}$.  Since
$ka^{-3/4}\le h^{-\xi+3\tau/4}$, summing these probabilities over
$k\le K$ is smaller than every inverse power of $h$.

If some $H_i\ge a$, fix its index and condition on the other pieces.
Proposition~\ref{c:height-atom}, followed by the heightwise bound of
Proposition~\ref{c:planar-atoms}, gives
\begin{align*}
 &\Pr\{S_k=w,\ H_i\ge a\}\\
 &\quad\le
 \sup_{n\ge a}p(H=n)
 \sum_m\sup_{H(v)=m}\Pr(S_{k-1}=v)
 \le a^{-7/4+o(1)}C(1+k)^{-4/3}.
\end{align*}
For $k=1$ the remaining sum is interpreted as the unit mass at zero.
Summing over the candidate index and then over $k\le K$ gives
\[
 C a^{-7/4+o(1)}K^{2/3}
 =h^{-5/4-2\xi/3+7\tau/4+o(1)}
 =o(h^{-5/4}).
\]
This is the step that needs both the one-piece height atom and the
heightwise supremum estimate.

For $k>h^{3/4+\xi}$, height discount alone gives
\[
 \sum_{k>h^{3/4+\xi}}\Pr\Bigl\{\sum_{i\le k}H_i=h\Bigr\}
 \le e\sum_{k>h^{3/4+\xi}}(\mathbb E e^{-H/h})^k
 \le Ch^{3/4}e^{-ch^\xi}.
\]
We are left with
$h^{3/4-\xi}<k\le h^{3/4+\xi}$.

Take $\xi$ small relative to $\epsilon$.  The length deficit in
\eqref{c:tail-atom-inputs}, with
$u=h^{-4/3+\epsilon}$, shows uniformly over these $k$ that
\[
 \Pr\{T_k<h^{4/3-\epsilon}\}
 \le e\,(\mathbb E e^{-uL})^k
 \le\exp\{-h^{9\epsilon/16-\xi-o(1)}\}.
\]
Summing in $k$ is negligible.

For the upper tail, split the pieces into two independent batches with
sizes comparable to $k$.  If $T_k>h^{4/3+\epsilon}$, one batch has
length exceeding $t=\tfrac12h^{4/3+\epsilon}$.  The length tail and
truncation yield, for any fixed small $\eta>0$,
\[
 \Pr\{\hbox{that batch has length}>t\}
 \le C_\eta k t^{-9/16+\eta}
 \le h^{-\delta}
\]
for some $\delta>0$, provided $\xi,\eta$ are sufficiently small
relative to $\epsilon$.  On the other, independent batch impose the
planar atom at the displacement needed to reach $w$.  This costs at
most $Ck^{-8/3}$.  Thus the sum of long-path masses over the remaining
$k$ is bounded by
\[
 C h^{-\delta}
 \sum_{k>h^{3/4-\xi}}k^{-8/3}
 \le C h^{-5/4+5\xi/3-\delta}=o(h^{-5/4}),
\]
after decreasing $\xi$ once more.  This use of separate batches avoids
any independence assertion within a piece or after endpoint
conditioning.

Both length exceptions therefore have absolute strip mass
$o(h^{-5/4})$.  Divide by the connector lower bound.  Restricting to a
finite domain can only decrease the exceptional mass, and its own
connector lower bound has the same order.  This proves both assertions.
\end{proof}
